\documentclass[11pt]{report}
\input{../pri-end/T0_preamble_standalone_A4_De}
% T0_preamble_patches.tex — LOKALER PLATZHALTER
% Im Repo durch die offizielle Version ersetzen.


\usepackage{longtable}
\usepackage{booktabs}

\hypersetup{
	pdftitle={250 · FFGFT — Information und schwarze Löcher},
	pdfauthor={Johann Pascher},
	pdfsubject={Grundinformation, topologische Invarianten, Hawking-Paradoxon, T4-Geometrie}
}

\fancyhead[L]{\small\textcolor{t0gray}{FFGFT · Information und schwarze Löcher}}
\fancyhead[R]{\small\textcolor{t0gray}{Johann Pascher · 2026}}

\title{%
	{\LARGE\bfseries\color{t0blue}
	 FFGFT --- Information und schwarze Löcher}\\[0.5em]
	{\large Was ist Grundinformation? Ontologie vs. Epistemologie}\\[1em]
	{\normalsize Johann Pascher \quad $\cdot$ \quad Mai 2026}\\[0.5em]
	{\small \url{https://github.com/jpascher/T0-Time-Mass-Duality}}\\[0.3em]
	{\small Zenodo: \href{https://doi.org/10.5281/zenodo.20117635}{10.5281/zenodo.20117635} (v1.1.0)}
}
\date{Mai 2026}

\begin{document}
	\maketitle

	% ==============================================================================
% 250_FFGFT_Schwarze_Loecher_Information_De_ch.tex
% Kapitelinhalt – einzubinden via Wrapper
% Autor: Johann Pascher · Mai 2026
% Überarbeitet nach DeepSeek-Review: alle 6 Verbesserungsvorschläge integriert
% ==============================================================================

\section*{Abstract}

Dieses Dokument klärt, was in FFGFT unter \emph{Information} im Zusammenhang mit schwarzen Löchern zu verstehen ist. Es zeigt, dass das Hawking-Informationsparadoxon auf einer Kategorienverwechslung beruht: Die klassische Formulierung behandelt Information als epistemisches Konzept. FFGFT unterscheidet davon die \textbf{Grundinformation}: die topologische Struktur der Windungszahlen auf $T^4$ --- vergleichbar mit einer topologischen Ladung, nicht mit einer Menge an Bits. Diese topologische Ladung bleibt \textbf{topologisch erhalten}, unabhängig von jedem Beobachter. Weiterhin wird gezeigt, dass die Kompaktifizierung der Zeitdimension auf $T^4$ eine der drei Voraussetzungen des Paradoxons untergräbt, und dass thermodynamische Irreversibilität mit topologischer Informationserhaltung vollständig verträglich ist.

\tableofcontents
\newpage

% ------------------------------------------------------------------------------
\section{Das Hawking-Informationsparadoxon und seine Voraussetzungen}

Stephen Hawking zeigte 1974, dass schwarze Löcher thermische Strahlung emittieren \cite{Hawking1974}. Diese Strahlung trägt keine Information über den spezifischen Zustand der eingefallenen Materie --- sie ist rein thermisch charakterisiert. Wenn das schwarze Loch vollständig verdampft, wäre die ursprüngliche Quanteninformation vernichtet. Das widerspricht der Unitarität der Quantenmechanik.

Das Paradoxon setzt drei Voraussetzungen voraus, die in FFGFT nicht alle gelten:

\begin{enumerate}
  \item Information ist die vollständige Beschreibung des Quantenzustands aller eingefallenen Teilchen.
  \item Zeit ist kontinuierlich und eindimensional gerichtet --- was ``in der Vergangenheit'' ins schwarze Loch gefallen ist, ist ``für immer'' unzugänglich.
  \item Die klassische Singularität ist physikalisch real --- der Zustand der Materie wird dort vernichtet.
\end{enumerate}

\begin{wichtig}[Warum entsteht das Paradoxon überhaupt?]
Hawking setzt (implizit) voraus, dass Information \emph{epistemisch} ist --- also das, was ein Beobachter rekonstruieren kann --- und dass Zeit kontinuierlich gerichtet ist. Beides ist in FFGFT nicht fundamental: Information ist eine topologische Ladung, keine epistemische Größe; und Zeit ist auf $T^4$ kompaktifiziert, nicht kontinuierlich. Das Paradoxon löst sich nicht durch Hinzufügen neuer Physik, sondern durch Präzisierung der Voraussetzungen.
\end{wichtig}

FFGFT weist alle drei Voraussetzungen in spezifischer Weise zurück --- nicht durch Behauptung, sondern durch die geometrische Struktur des Frameworks.

% ------------------------------------------------------------------------------
\section{Was Information in FFGFT \emph{nicht} ist}

\subsection{Nicht unsere Beschreibung des schwarzen Lochs}

Information im FFGFT-Sinne ist \textbf{nicht} unser Wissen oder Verständnis des schwarzen Lochs. Das Hawking-Paradoxon ist ein \textbf{epistemisches} Problem: Es fragt, ob ein Beobachter außerhalb des Horizonts den ursprünglichen Quantenzustand vollständig rekonstruieren kann. FFGFT beantwortet eine andere Frage: Existiert die fundamentale topologische Struktur weiterhin?

\begin{wichtig}
Information im Sinne von FFGFT ist keine epistemische Kategorie (was wir wissen können), sondern eine \textbf{ontologische} Kategorie (was strukturell existiert). Die Unfähigkeit eines äußeren Beobachters, Information zu rekonstruieren, ist ein Zugangsproblem --- kein Existenzproblem.

\textbf{Messbare Konsequenz:} Wäre die topologische Information wirklich vernichtet, wäre die Hawking-Strahlung exakt schwarzkörperartig. Das modifizierte Spektrum (durch die $T^4$-Windungsstruktur) ist die falsifizierbare Differenz --- damit wird das Zugangsproblem empirisch geerdet.
\end{wichtig}

\subsection{Nicht die Summe aller Detailzustände}

Information im FFGFT-Sinne ist auch nicht die Summe aller Informationen, die zur vollständigen Beschreibung aller Teilchen nötig wären. Diese Detailzustände können durch Thermalisierung epistemisch verloren gehen. Das ist thermodynamische Irreversibilität, kein ontologischer Informationsverlust.

% ------------------------------------------------------------------------------
\section{Was Information in FFGFT \emph{ist}: Die Grundinformation}

\subsection{Topologische Invarianten als Grundinformation}

In FFGFT sind Teilchen Windungskonfigurationen auf dem 4D-Identifikationstorus $T^4$ mit Quantenzahlen $(n_\theta, n_\varphi, n_3, n_4)$. Diese Windungszahlen sind \textbf{topologische Invarianten}.

\begin{keyresult}[Definition der Grundinformation]
Die Grundinformation eines physikalischen Systems in FFGFT ist die \textbf{topologische Ladung} seiner Bestandteile: die Windungszahlen $(n_\theta, n_\varphi, n_3, n_4)$ auf $T^4$. Diese definieren, \emph{welche Art} von Teilchen existieren --- nicht \emph{wo} sie sind, nicht \emph{wie} sie sich verhalten, nicht \emph{in welchem genauen Quantenzustand} sie sich befinden.

Grundinformation ist keine \emph{Menge an Bits} (Shannon), sondern eine \textbf{topologische Ladung} --- vergleichbar mit einer magnetischen Monopolladung oder einer Windungszahl in einer Stringtheorie. Das entkoppelt FFGFT sofort von der üblichen Bit-orientierten Informationsparadoxon-Debatte.
\end{keyresult}

Topologische Invarianten können nicht durch kontinuierliche Transformationen auf null reduziert werden. Ein Zustand mit Windungszahl $n = 1$ kann nicht kontinuierlich in $n = 0$ überführt werden, ohne eine topologische Diskontinuität --- mit messbaren physikalischen Konsequenzen.

\begin{erkenntnis}
\textbf{Kein Vernichtungsoperator wirkt auf topologische Invarianten.} Das ist der zentrale Satz: Es gibt keinen physikalischen Prozess, der Windungszahlen auf null reduzieren kann, ohne eine messbare Diskontinuität zu erzeugen. Singularitäten im klassischen Sinne wären ein solcher Prozess --- deshalb sind sie in FFGFT strukturell ausgeschlossen.
\end{erkenntnis}

\subsection{Die Grundformel als Strukturerhaltung}

Die fundamentale Relation $\tilde{T} \cdot m = 1$ verbindet T-Feld und Masse an jedem Punkt. Im Grenzfall extremer Kompression:
\begin{equation}
  \tilde{T} \rightarrow L_0/c = \xipar \cdot t_P > 0
\end{equation}
Das T-Feld erreicht nie null --- der Minimalwert ist $L_0/c$. Die Windungsstruktur der Feldkonfiguration bleibt erhalten.

\subsection{Algebraischer Schutz vor Singularitäten (Dok.~078)}

Hier liegt der eleganteste Beweis --- nicht durch ein Postulat, sondern durch die algebraische Struktur der Grundrelation selbst.

Aus $\tilde{T} \cdot m = 1$ folgt:
\begin{equation}
  \tilde{T} = \frac{1}{m}
\end{equation}

Wenn $m \rightarrow \infty$, dann $\tilde{T} \rightarrow 0$. Aber $\tilde{T} = 0$ exakt zu erreichen würde $m = \infty$ exakt erfordern --- algebraisch unmöglich.

\begin{keyresult}[Algebraischer Schutz --- Dok.~078]
$\tilde{T} \cdot m = 1$ wirkt wie ein algebraischer Schutz vor Singularitäten:
\begin{itemize}
  \item $m \rightarrow \infty$ impliziert $\tilde{T} \rightarrow 0$ --- aber nur \textbf{asymptotisch}, nie exakt.
  \item $\tilde{T} = 0$ exakt würde $m = \infty$ exakt erfordern --- algebraisch unmöglich.
  \item Die Singularität ist kein verbotener Zustand, sondern ein \textbf{unerreichbarer Grenzwert}.
\end{itemize}
Das ist kein geometrisches und kein quantenmechanisches Argument --- es ist ein \textbf{algebraisches Argument} aus der Grundrelation allein. Es gilt unabhängig von der Größe des schwarzen Lochs und unabhängig davon wie weit man die Kompression treibt.
\end{keyresult}

\subsection{Der Interior-Operator: Was bei $L_0$ passiert}

Was beim Grenzwert $L_0$ im Inneren eines schwarzen Lochs geschieht, ist in FFGFT nicht vollständig offen.

\textbf{Das T-Feld am Minimum.}
$\tilde{T}_\text{min} = \xipar \cdot t_P > 0$ --- ein endlicher, positiver, nicht-singulärer Wert.

\textbf{Maximale Windungsdichte.}
Ein Windungsquant pro $L_0^4$ Raumzeitvolumen --- analogt zur unteren Schranke in einem Gitter mit endlichem Volumen. Das ist die dichteste stabile Konfiguration auf der $\xipar$-Leiter (Niveau $N = 0$). Unterhalb endet der Zustandsraum.

\textbf{Lagrange-Steifigkeit.}
Der Term $-\frac{1}{2}m_T^2(\partial m)^2$ mit $m_T = \xipar/\ell$ (Dok.~201) verhindert, dass $\partial m$ divergiert --- algebraische Eigenschaft der Feldgleichungen, keine externe Kraft.

\textbf{Topologische Bounce-Bedingung.}
Keine stabilen Konfigurationen unterhalb $N = 0$ $\Rightarrow$ keine Zustände in die das System übergehen kann. Das ist eine \textbf{Zustandsraumgrenze}, kein Kraft-Gegenkraft-Mechanismus.

\textbf{Zeit im Inneren.}
$\tilde{T} \rightarrow L_0/c$: Lokaler Zeitquant am Minimum. Zeit friert nicht ein ($\tilde{T} = 0$ wäre nötig), aber die lokale Zeitstruktur erreicht ihren kleinsten möglichen Wert. Das modifizierte Hawking-Spektrum entsteht genau an dieser Grenzfläche.

\textbf{Formales Kriterium für den Übergangs-Operator.}
Ein Interior-Operator müsste die Abbildung einer eingehenden Windungskonfiguration $(n_i)$ auf die extremale Konfiguration bei $L_0$ beschreiben --- unter Erhaltung der topologischen Ladung, aber mit maximaler Windungsdichte. Die $T^4$-Algebra legt nahe, dass dieser Operator \textbf{unitär} ist, aber nichttrivial mit dem thermodynamischen Limes vertauscht. Das ist die offene Forschungsfrage --- präzise benannt, keine prinzipielle Lücke.

\begin{erkenntnis}[Inneres eines schwarzen Lochs in FFGFT]
Das Innere ist kein unbekanntes Territorium, sondern ein Extremzustand:
\begin{itemize}
  \item $\tilde{T}_\text{min} = \xipar \cdot t_P > 0$ --- endlich, nicht singulär
  \item Maximale Windungsdichte: $1/L_0^4$ pro Raumzeitvolumen
  \item Lagrange-Steifigkeit: $\partial m$ divergiert nicht
  \item Zustandsraumgrenze bei $N = 0$
  \item Unitärer Übergangs-Operator (Details offen)
\end{itemize}
\end{erkenntnis}

\subsection{Grundinformation ist nicht die Gesamtbeschreibung}

Die Grundinformation ist die \textbf{Klasse} der Konfigurationen (topologische Ladung), nicht ihre \textbf{Instanz} (Quantenzustand):
\begin{itemize}
  \item Klasse: ``Es gibt ein Elektron'' (Windungszahlen definiert)
  \item Instanz: ``Das Elektron ist an Position $x$ mit Spin $+\frac{1}{2}$'' (Quantenzustand)
\end{itemize}
FFGFT erhält die Klasse. Die Instanz kann thermodynamisch irreversibel werden.

% ------------------------------------------------------------------------------
\section{Schwarze Löcher als Inversion --- nicht als Vernichtung}

\subsection{Die topologische Struktur bleibt invertiert}

Ein schwarzes Loch ist kein Vernichtungszustand, sondern eine \textbf{Inversion} der normalen Materiestruktur: Windungskonfigurationen maximal komprimiert, T-Feld am Minimum $L_0/c$. Was als Elektron einfiel, ist weiterhin topologisch als Elektron klassifizierbar --- in einem extremen Feldregime.

\subsection{Irreversibilität ist thermodynamisch, nicht ontologisch}

\begin{erkenntnis}
Irreversibilität und Informationserhaltung operieren auf verschiedenen Ebenen:
\begin{itemize}
  \item \textbf{Thermodynamisch:} Gerichtet, irreversibel, Entropie nimmt zu.
  \item \textbf{Topologisch:} Windungszahlen \textbf{topologisch erhalten}. Kein Vernichtungsoperator wirkt auf topologische Invarianten.
\end{itemize}
Das Hawking-Paradoxon entsteht durch Verwechslung dieser beiden Ebenen.
\end{erkenntnis}

% ------------------------------------------------------------------------------
\section{Der Strahlungsmechanismus: Knotentunneln und modifiziertes Spektrum}

Die vorangehenden Abschnitte haben geklärt, was mit der \emph{Information} geschieht. Dieser Abschnitt klärt den \emph{Mechanismus} der Strahlung selbst --- warum ein schwarzes Loch überhaupt emittiert --- und leitet die FFGFT-spezifische Modifikation des Spektrums aus erster Linie her.

\subsection{Temperatur und Mechanismus}

Die Hawking-Temperatur bleibt in FFGFT \textbf{unverändert} (Dok.~201, §7.4):
\begin{equation}
  T_H = \frac{\hbar \kappa}{2\pi k_B} = \frac{\hbar c^3}{8\pi G M k_B},
  \qquad \kappa = \frac{c^4}{4GM}.
\end{equation}
Die Oberflächengravitation $\kappa$ ist geometrisch festgelegt und wird durch die $T^4$-Struktur nicht modifiziert. Was sich ändert, ist der \textbf{Emissionsmechanismus}: statt Teilchen-Antiteilchen-Paarbildung in einem kontinuierlichen Vakuum emittiert FFGFT \textbf{inkohärente Knotenanregungen}, die durch die Phasenbarriere am Horizont \textbf{tunneln} (Dok.~201, §7.4; Dok.~049). Dies ist die FFGFT-Übersetzung des Parikh-Wilczek-Tunnelbildes: die Quanten sind Windungsanregungen des $T$-Feldes, nicht Moden eines kontinuierlichen Feldes.

Numerisch (Sonnenmasse, $M = M_\odot$): $T_H \approx 6{,}17 \times 10^{-8}$ K; Schwarzschild-Radius $r_s \approx 2{,}95$ km. Da $T_H \propto 1/M$, beschleunigt sich die Verdampfung mit abnehmender Masse.

\subsection{Entropie als Knotenzahl}

Die Bekenstein-Hawking-Entropie ist in FFGFT die \textbf{Konfigurationsentropie der Knoten}:
\begin{equation}
  S = \frac{A}{4 \ell_P^2} = N_\text{Knoten} \cdot \ln 2,
  \qquad N_\text{Knoten} \propto \frac{A}{\xipar^2}.
\end{equation}
Jeder Knoten trägt genau ein Bit ($\ln 2$). Für $M = M_\odot$: $S/k_B \approx 1{,}05 \times 10^{77}$, also $N_\text{Knoten} \approx 1{,}5 \times 10^{77}$. Das reproduziert das Flächengesetz exakt --- ohne freie Anpassung.

\subsection{Herleitung der Spektrum-Modifikation aus der Gitterdispersion}

Die zentrale FFGFT-Vorhersage --- ein vom reinen Schwarzkörper abweichendes Spektrum --- folgt nicht aus einem Ansatz, sondern aus der \textbf{minimalen Länge} $L_0 = \xipar \cdot \ell_P$. Eine minimale Länge impliziert eine diskrete Gitterstruktur, und ein Gitter besitzt eine modifizierte Dispersionsrelation (analog zur Phonon-Dispersion im Festkörper, Debye-Theorie):
\begin{equation}
  E(k) = \frac{2\hbar c}{L_0} \sin\!\left(\frac{k L_0}{2}\right)
       = \hbar c k \left[1 - \frac{(k L_0)^2}{24} + \mathcal{O}\!\left((kL_0)^4\right)\right].
\end{equation}
Die führende Korrektur ist damit:
\begin{equation}
  \frac{\Delta E}{E} = -\frac{(k L_0)^2}{24}
  = -\frac{1}{24}\left(\frac{E}{E_\text{max}}\right)^2,
  \qquad E_\text{max} = \frac{\hbar c}{L_0} = \frac{E_P}{\xipar}.
\end{equation}
Der Cutoff $E_\text{max} = E_P/\xipar \approx 7500 \cdot E_P \approx 9{,}2 \times 10^{22}$ GeV liegt oberhalb der Planck-Energie. In Einheiten der Hawking-Temperatur ($E = x \, k_B T_H$):
\begin{equation}
  \frac{\Delta E}{E} = -\frac{x^2}{24}\left(\frac{k_B T_H}{E_\text{max}}\right)^2.
\end{equation}

\subsection{Größenordnung und Falsifizierbarkeit --- eine ehrliche Einordnung}

Diese Herleitung hat das korrekte Grenzverhalten, aber ihre praktische Konsequenz muss \textbf{ehrlich genannt} werden: für reale schwarze Löcher ist die Modifikation unmessbar klein.

Für ein stellares schwarzes Loch ($M = M_\odot$) ist die typische Quantenenergie $k_B T_H$ winzig gegenüber $E_\text{max}$:
\begin{equation}
  \frac{k_B T_H}{E_\text{max}} \approx 5{,}8 \times 10^{-44}
  \quad\Rightarrow\quad
  \frac{\Delta E}{E}\bigg|_{x=10} \approx -1{,}4 \times 10^{-86}.
\end{equation}
Die Modifikation wird erst signifikant, wenn $k_B T_H$ sich $E_\text{max}$ nähert --- also wenn das schwarze Loch bis zur Massenskala
\begin{equation}
  M_\text{krit} = \frac{c^2 L_0}{8\pi G} \approx 1{,}15 \times 10^{-13} \text{ kg}
\end{equation}
verdampft ist. Das ist genau die \textbf{Remnant-Skala}. Die Spektrum-Modifikation entsteht damit exakt an der $L_0$-Grenzfläche und am Verdampfungs-Endpunkt --- konsistent mit dem Interior-Operator-Abschnitt.

\begin{erkenntnis}[Korrektur eines naiven Ansatzes]
Ein naiver Modulationsansatz der Form $\exp(-\xipar \, x^2 / D_f)$ ergibt künstlich große Abweichungen ($\sim 0{,}4\%$), die physikalisch \textbf{nicht existieren}. Die korrekte Abweichung, aus der Gitterdispersion hergeleitet, skaliert wie $(k_B T_H/E_\text{max})^2$ und beträgt $\sim 10^{-86}$ für reale schwarze Löcher. Eine Theorie, die eine messbare $0{,}4\%$-Abweichung behauptete, wäre bereits falsifiziert. Die korrekte FFGFT-Vorhersage ist die viel schwächere, aber konsistente Aussage: das Spektrum ist nur am Verdampfungs-Endpunkt modifiziert.
\end{erkenntnis}

\subsection{Verdampfungs-Endpunkt: Remnant statt Vernichtung}

Während das schwarze Loch verdampft, schrumpft $M$ und mit ihm der Kernradius. Nahe $M_\text{krit}$ erreicht das System die Zustandsraumgrenze $N = 0$ (maximale Windungsdichte $1/L_0^4$): es gibt keine dichteren stabilen Konfigurationen, in die es übergehen könnte. Das Ergebnis ist ein \textbf{stabiles Remnant} (Dok.~201, §7.5):
\begin{itemize}
  \item endliche Größe $\sim L_0$,
  \item endliche Temperatur,
  \item topologische Ladung im Remnant-Knotenzustand erhalten.
\end{itemize}
Die ursprüngliche Grundinformation bleibt damit \textbf{diesseits} kodiert --- im finalen stabilen Knotenzustand. Es gibt keinen Transfer auf eine ``andere Seite'' und keine Vernichtung. Die thermodynamische Zeitrichtung der Verdampfung bleibt irreversibel; die topologische Erhaltung ist davon unberührt.

% ------------------------------------------------------------------------------
\section{Die kompaktifizierte Zeitdimension}

\subsection{Zeit ist nicht kontinuierlich --- sie ist kompaktifiziert}

In FFGFT ist die Zeitdimension auf $T^4$ \textbf{kompaktifiziert} --- die Windungszahl $n_4$ ist eine topologische Invariante wie die anderen.

\begin{wichtig}
Kompaktifizierung bedeutet \textbf{Periodizität der Feldkonfigurationen} in der Zeitrichtung --- nicht Umkehrbarkeit der thermodynamischen Zeit. Ein makroskopischer Prozess bleibt irreversibel. Die $T^4$-Topologie ist zeitrichtungsneutral; die thermodynamische Zeitrichtung ist emergent, nicht fundamental.
\end{wichtig}

\subsection{Was die Kompaktifizierung ändert}

Die Aussage ``Information ist für immer in der Vergangenheit verloren'' setzt eine absolute ontologische Vergangenheitsgrenze voraus. In einer kompaktifizierten Zeitdimension ist diese Absolutheit nicht selbstverständlich: Die topologische Feldstruktur ist nicht auf die Zeitrichtung des Beobachters referenziert. $n_4$ bleibt topologisch erhalten.

\subsection{Erlebte Zeit ist eindeutig gerichtet}

Die erlebte Zeitrichtung ist ein thermodynamisches Phänomen (Entropiezunahme), kein geometrisches Primitiv. In FFGFT ist die Zeitrichtung emergent, nicht fundamental --- die Kompaktifizierung operiert auf der topologischen Ebene und überlastet die Theorie nicht mystisch.

% ------------------------------------------------------------------------------
\section{Was FFGFT konkret behauptet --- und was nicht}

\begin{enumerate}
  \item \textbf{Keine klassische Singularität:} $L_0 = \xipar \cdot \ell_P$ als absolute untere Dichtegrenze.
  \item \textbf{Topologische Grundinformation (Ladung) erhalten:} Windungszahlen als topologische Invarianten, algebraisch geschützt durch $\tilde{T} \cdot m = 1$ (Dok.~078).
  \item \textbf{Modifiziertes Hawking-Spektrum:} Falsifizierbare Vorhersage --- Abweichung vom Schwarzkörperspektrum, aus der Gitterdispersion bei minimaler Länge $L_0$ hergeleitet ($\Delta E/E = -(E/E_\text{max})^2/24$). Für reale schwarze Löcher unmessbar klein; signifikant erst an der Remnant-Skala.
  \item \textbf{Gravitationswellen-Echos:} Nach Verschmelzungen als weiterer Test.
\end{enumerate}

\textbf{Nicht behauptet:} Vollständige Zustandsrekonstruktion; Träger-Recovery; vollständig spezifizierter Übergangs-Operator nahe $N = 0$; Aufhebung thermodynamischer Irreversibilität.

\begin{keyresult}[Präzise FFGFT-Aussage]
FFGFT behauptet: Die \textbf{topologische Grundladung} --- die Windungszahlen der eingefallenen Materie --- ist \textbf{topologisch erhalten}. Sie ist nicht rekonstruierbar durch einen äußeren Beobachter (epistemisches Zugangsproblem), aber sie existiert weiterhin als geometrische Eigenschaft des $T^4$-Feldes.

Wäre sie wirklich vernichtet, wäre das Hawking-Spektrum exakt thermisch --- das modifizierte Spektrum ist die messbare Differenz.
\end{keyresult}

% ------------------------------------------------------------------------------
\section{Verbindung zur $\xipar$-Rekursion}

Die $\xipar$-Leiter definiert welche Windungskonfigurationen auf $T^4$ stabil sind. Schwarze Löcher repräsentieren Konfigurationen am unteren Ende ($N = 0$) --- innerhalb der $\xipar$-Struktur, nicht außerhalb. Auch die ``invertierte'' Konfiguration eines schwarzen Lochs ist Teil der $\xipar$-Leiter und topologisch wohlbestimmt.

% ------------------------------------------------------------------------------
\section*{Zusammenfassung}

\begin{longtable}{>{\bfseries}p{6cm} p{8cm}}
\toprule
Frage & FFGFT-Antwort \\
\midrule
Warum entsteht das Paradoxon überhaupt? & Hawking setzt Information als epistemisch und Zeit als kontinuierlich voraus --- beides in FFGFT nicht fundamental \\
Was ist Grundinformation in FFGFT? & Topologische Ladung: Windungszahlen $(n_\theta, n_\varphi, n_3, n_4)$ auf $T^4$ --- keine Bits \\
Was ist keine Information? & Unser Wissen über den Zustand; Träger-Identität; vollständige Zustandsbeschreibung \\
Geht Information verloren? & Topologische Ladung: nein (topologisch erhalten). Detailzustand: Zugangsproblem \\
Empirische Erdung des Zugangsproblems & Modifiziertes Hawking-Spektrum: wäre Information vernichtet, wäre Strahlung exakt thermisch \\
Algebraischer Schutz (Dok.~078) & $m \rightarrow \infty$ impliziert $\tilde{T} \rightarrow 0$ nur asymptotisch --- Singularität algebraisch unerreichbar \\
Was passiert bei $L_0$? & $\tilde{T}_\text{min} > 0$ · Maximale Windungsdichte · Lagrange-Steifigkeit · Zustandsraumgrenze $N=0$ \\
Gibt es eine Singularität? & Nein --- $L_0$ als Untergrenze; Zustandsraum endet bei $N = 0$ \\
Ist der Prozess irreversibel? & Thermodynamisch ja. Topologisch: Grundstruktur bleibt erhalten \\
Zeitkompaktifizierung & Periodizität der Feldkonfigurationen, nicht Umkehr der thermodynamischen Zeit \\
Falsifizierbare Tests & Hawking-Spektrum (nicht exakt thermisch); Gravitationswellen-Echos \\
Was bleibt offen? & Übergangs-Operator $(n_i) \rightarrow$ extremale Konfiguration bei $L_0$ --- unitär, Details offen \\
\bottomrule
\end{longtable}


	\bibliographystyle{plainnat}
	\bibliography{../bib/FFGFT_Epistemologie}

\end{document}
